\section{積分の微分}\label{節:変数上下端}
$\int_a^x f(t)\dd t$ の微分は,基本定理で分かった.では,\,$\int_x^{x^2}\e^{t^2}\dd t$ のように,上端や下端が $x$ の関数になっている場合はどうだろうか.下端を定数にした形に書き直し,合成関数の微分に帰着させよう.

$f$ がある区間で連続で,\,$u(x),v(x)$ が微分可能,その値がこの区間内にあるとする.区間内の点 $c$ を固定し,
\[
 A(s)=\int_c^s f(t)\dd t
\]
と置く.区間の加法性より
\begin{align*}
 \int_{u(x)}^{v(x)}f(t)\dd t
 &=\int_c^{v(x)}f(t)\dd t-\int_c^{u(x)}f(t)\dd t\\
 &=A(v(x))-A(u(x))
\end{align*}
である.両辺を $x$ で微分し,合成関数の微分と $A'=f$ を使うと,
\[
 \boxed{\frac{\dd}{\dd x}\int_{u(x)}^{v(x)}f(t)\dd t
 =f(v(x))v'(x)-f(u(x))u'(x)}
\]
を得る.下端からの寄与に負号が付くのは,区間の加法性でその部分を引いているからである.\,$u(x)$ と $v(x)$ の大小が途中で入れ替わっても,向きを付けた積分の約束によりこの式を使える.

\begin{bookpoint}
\textbf{POINT}\par
上端 $v(x)$ を $f$ に代入して $v'(x)$ を掛け,下端 $u(x)$ を代入して $u'(x)$ を掛けたものを引く.
\end{bookpoint}

\begin{example}{積分の微分}
$G(x)=\int_x^{x^2}\e^{t^2}\dd t$ を微分せよ.
\end{example}
\paragraphbf{方針}
上端 $x^2$,下端 $x$ を代入し,それぞれの微分を掛けて差をとる.

\paragraphbf{解答}
$G'(x)=2x\e^{x^4}-\e^{x^2}$

\paragraphbf{解法}
被積分関数 $\e^{t^2}$ はすべての実数で連続である.上端 $x^2$ の微分は $2x$,下端 $x$ の微分は $1$ だから,
\begin{align*}
 G'(x)&=\e^{(x^2)^2}\cdot2x-\e^{x^2}\cdot1\\
 &=2x\e^{x^4}-\e^{x^2}
\end{align*}
となる.上端を $t$ に代入するので,最初の指数は $x^4$ になる.\,$\e^{t^2}$ の原始関数を具体的な式で表す必要はない.

\paragraphbf{被積分関数に $x$ を含む場合}
この公式は,被積分関数 $f(t)$ に外側の変数 $x$ が含まれない場合のものである.たとえば
\[
 H(x)=\int_0^x xt\dd t
\]
では,積分の間は $x$ を定数として外に出せて,
\begin{align*}
 H(x)&=x\int_0^x t\dd t\\
 &=x\cdot\frac{x^2}{2}\\
 &=\frac{x^3}{2}
\end{align*}
となる.したがって $H'(x)=\frac{3x^2}{2}$ である.公式のように上端を代入しただけの $x\cdot x=x^2$ では足りない.外に出した因子 $x$ も,積分全体を $x$ で微分するときには変わるからである.実際,積の微分を使えば
\[
 H'(x)=\int_0^x t\dd t+x\cdot x=\frac{x^2}{2}+x^2
\]
となり,同じ結果を得る.
